% SGA 3，第 VII_B 篇，英文重建稿——第 1 轮。
% 范围：权威底本分篇第 31--34 页，第 1.3.6 节全文，止于第 1.4 节之前。
% 合并阅读本页码： 544--547.
% 再版印刷页码： 533--536.
% 本范围内原版页边定位： 504--506.
% 权威底本： Polo--Gille Exp7B-13oct24.pdf,
% SHA-256 D303690E0195654BBE7C3DAA5D89EA23E24CD709A3575217234AC19AE8836694.
% 曾参考 Jacob Reinhold 的 CC BY 4.0 Markdown 作为比对谱系。
% 原英文 Loop 1 阶段说明：唯一的图以权威 PDF 的精确裁片保留，其余陈列数学以原生 LaTeX 重建。当前中文工程已采用原生图。
\subsubsection*{1.3.6}\phantomsection \label{sga3:VIIB:1.3.6}\addcontentsline{toc}{subsubsection}{\SGARefLink{sga3:VIIB:1.3.6}{1.3.6}. 带滤过的平坦余代数}

在本讲余下的部分，我们将多次通过给出余代数 \(\mathbf H(X)\) 来定义拓扑平坦的形式 \(k\)-簇 \(X\)。这时需要借助 \(\mathbf H(X)\) 解释 \(X\) 的某些几何性质。例如，设给定结构态射 \(X\to\operatorname{Spf}(k)\) 的一个截面 \(\sigma\)，问在什么条件下 \(\sigma\) 在底拓扑空间之间诱导同构。\footnote{\label{sga3:VIIB:nde:62}在下面的 \SGARefLink{sga3:VIIB:1.3.6.A}{引理 1.3.6.A} 中，编者补充了 \textup{\SGARefLink{sga3:vii:paragraph:5-1:equation:i}{(i)}} 和 \textup{\SGARefLink{sga3:vii:paragraph:5-1:equation:ii}{(ii)}} 的详细证明，并增添了 \textup{\SGARefLink{sga3:VIIB:1.3.6.A:item:iii}{(iii)}}。}

\textit{首先假设 \(k\) 是 Artin 环。}设 \((H,\Delta,\varepsilon)\) 是平坦的 \(k\)-余代数，置 \(H^+=\ker(\varepsilon)\)，并设 \(A=H^*\) 为 \(H\) 的对偶预有限 \(k\)-代数。设给定一个 \(k\)-余代数态射 \(k\to H\)，也就是说，给定满足下式的元素 \(\phi\in H\)：\[
  \varepsilon(\phi)=1,
  \qquad
  \Delta(\phi)=\phi\otimes\phi.
\] 一方面，\(\phi\) 定义连续代数同态 \(\Phi:A\to k\)，从而定义投影 \(\operatorname{Spf}(A)\to\operatorname{Spf}(k)\) 的截面 \[
  \sigma:\operatorname{Spf}(k)\longrightarrow\operatorname{Spf}(A)
\]。

另一方面，令 \(H_0=k\phi\)，并对 \(n\geq1\) 令 \[
  H_n
  =
  \{x\in H\mid
    \Delta(x)-x\otimes\phi\in H_{n-1}\otimes H^+\}.
\]，以此定义 \(H\) 的 \(k\)-子模。若置 \(H_{-1}=0\)，则此式也适用于 \(n=0\)。对 \(n\) 归纳可知 \(H_{n-1}\subset H_n\)。我们称 \[
  H_0\subset H_1\subset\cdots
\] 为\emph{由 \(\phi\) 定义的 \(H\) 的滤过}。

\medskip\noindent\textbf{注。}\phantomsection \label{sga3:VIIB:1.3.6.remark} 由 \[
  \Delta(H_n)
  \subset
  H_n\otimes H_0\oplus H_{n-1}\otimes H^+,
\] 可得 \(\Delta(H_n)\subset H_n\otimes H\)。由于 \(\Delta\) 余交换，即 \(\tau\circ\Delta=\Delta\)，其中 \(\tau(a\otimes b)=b\otimes a\)，同样有 \(\Delta(H_n)\subset H\otimes H_n\)。当 \(H/H_n\) 在 \(k\) 上平坦时，便知 \(H_n\) 是 \(H\) 的子余代数 \textup{（另见下面的 \SGARefLink{sga3:VIIB:1.3.6.A}{1.3.6.A}\,\textup{\SGARefLink{sga3:VIIB:1.3.6.A:item:iii}{(iii)}}）}。但一般而言，映射 \[
  \Delta:H_n\longrightarrow H_n\otimes H
\] 并不通过 \(H_n\otimes H_n\) 分解。\footnote{\label{sga3:VIIB:nde:63}例如，设 \(k_0\) 为域，令 \(k=k_0[T]/(T^n)\)，其中 \(n\geq4\)，并令 \(H=k\phi\oplus kx\)，且 \(\varepsilon(x)=0\)，\[
  \Delta(x)
  =
  x\otimes\phi+\phi\otimes x+t\,x\otimes x,
\]，其中 \(t\) 是 \(T\) 在 \(k\) 中的像。则对 \(i=0,\ldots,n\) 有 \(H_i=k\phi\oplus t^{\,n-i}x\)，但当 \(2\leq i\leq n-2\) 时，元素 \(\Delta(t^{\,n-i}x)\) 不属于 \(H_i\otimes H_i\) 在 \(H\otimes H\) 中的像。}

\medskip\noindent\textbf{引理 1.3.6.A。}\phantomsection \label{sga3:VIIB:1.3.6.A} {\itshape 设 \(k\) 为 Artin 环，\(H\) 为平坦 \(k\)-余代数，\(A=H^*\) 为它的对偶预有限 \(k\)-代数。设 \(\phi\in H\) 满足 \(\varepsilon(\phi)=1\) 和 \(\Delta(\phi)=\phi\otimes\phi\)。设 \(\Phi:A\to k\) 为相应的连续代数同态，\(\sigma:\operatorname{Spf}(k)\to\operatorname{Spf}(A)\) 为 \(\operatorname{Spf}(A)\to\operatorname{Spf}(k)\) 的相应截面，\((H_n)\) 为由 \(\phi\) 定义的 \(H\) 的滤过。置 \(I=\ker\Phi\)。\begin{enumerate}\item[\textup{(i)}]对每个 \(n\geq1\)，\(H_{n-1}\) 是 \(I^n\) 的闭包 \(\overline{I^n}\) 在 \(H\) 中的正交补。

\item[\textup{(ii)}]因此，\(\sigma\) 在底集合之间诱导双射，当且仅当 \[
    H=\bigcup_n H_n.
  \]。\footnote{\label{sga3:VIIB:nde:64}此时称余代数 \(H\) 为\emph{连通的}；参见增补 \SGARefLink{sga3:VIIB:2.9}{2.9}。}

\item[\textup{\SGAInlineRefTarget{sga3:VIIB:1.3.6.A:item:iii}(iii)}]若再假设每个 \(H/H_n\) 在 \(k\) 上平坦，则对每个 \(n\geq0\) 有 \begin{equation}
    \Delta(H_n)
    \subset
    \sum_{i=0}^n H_i\otimes H_{n-i}.
    \tag{*}
    \label{sga3:VIIB:1.3.6.A:star}
  \end{equation}。特别地，此时每个 \(H_n\) 都是 \(H\) 的子余代数。\end{enumerate}}

\textit{证明。}对每个 \(x\in H\)，映射 \[
  A\longrightarrow k,\qquad f\longmapsto f(x),
\] 连续。因此，若 \(I^n\) 被 \(x\) 零化，则其闭包 \(\overline{I^n}\) 也被零化。对 \(n\geq1\)，置 \[
  (\overline{I^n})^\perp
  =
  \{x\in H\mid
    f(x)=0\text{ 对所有 }f\in\overline{I^n}\}.
\]。

先设 \[
  \sigma:\mathfrak m\longmapsto\Phi^{-1}(\mathfrak m)
\] 是从 \(\operatorname{Spf}(k)\) 到 \(\operatorname{Spf}(A)\) 的双射。由于 \(I\) 包含于各理想 \(\Phi^{-1}(\mathfrak m)\) 的交，由 \SGARefLink{sga3:VIIB:0.1.2}{0.1.2} 可知理想序列 \((\overline{I^n})\) 趋于零。设 \(x\in H\)。子模 \[
  J(x)=\{f\in A\mid f(x)=0\}
\] 在 \(A\) 中既开又闭，故对所有充分大的 \(n\)，它包含 \(\overline{I^n}\)。换言之，对所有充分大的 \(n\)，有 \(x\in(\overline{I^n})^\perp\)。

反之，设 \[
  H=\bigcup_n(\overline{I^n})^\perp.
\]。设 \(\mathfrak p\) 是 \(A\) 的开素理想。根据 \(A\) 上拓扑的定义 \textup{（\SGARefLink{sga3:VIIB:0.2.2}{0.2.2}）}，\(\mathfrak p\) 包含形如 \[
  \mathscr V(x_1,\ldots,x_s)
  =
  \{f\in A\mid f(x_1)=\cdots=f(x_s)=0\}.
\] 的开 \(k\)-子模。由假设，存在整数 \(n\) 使得 \[
  x_1,\ldots,x_s\in(\overline{I^n})^\perp;
\]，因此 \(\overline{I^n}\subset\mathfrak p\)。此外，由于 \(k\) 是 Artin 环，\[
  \operatorname{Spf}(k)=\{\mathfrak m_1,\ldots,\mathfrak m_r\}
\] 是有限集，且存在整数 \(t\geq1\) 使得 \[
  \left(\prod_i\mathfrak m_i\right)^t=0.
\]。由此得 \[
  \prod_i\Phi^{-1}(\mathfrak m_i)^t\subset I.
\]。所以 \(\mathfrak p\) 包含各理想 \(\Phi^{-1}(\mathfrak m_i)^{tn}\) 的乘积。由于 \(\mathfrak p\) 为素理想，它包含某个 \(\Phi^{-1}(\mathfrak m_i)\)，因而与之相等。于是我们证明了 \footnotemark \begin{equation}
  \sigma\text{ 是双射}
  \quad\Longleftrightarrow\quad
  H=\bigcup_n(\overline{I^n})^\perp.
  \tag{65}
  \label{sga3:VIIB:1.3.6:eq65}
\end{equation}\footnotetext{\label{sga3:VIIB:nde:65}不假设 \(H\) 余交换而仍假设 \(k\) 为交换 Artin 环时，这仍然成立。此时，\(A=H^*\) 的零点邻域基由使 \(A/J\) 在 \(k\) 上具有有限长度的双边理想 \(J\) 给出，而上述证明表明，\(H=\bigcup_n(\overline{I^n})^\perp\) 当且仅当各 \(\Phi^{-1}(\mathfrak m_i)\) 是 \(A\) 的全部开素理想。}

又有 \(H=k\phi\oplus H^+\)；设 \(\pi:H\to H^+\) 为以 \(k\phi\) 为核的投影。对 \(n\geq0\)，设 \[
  \Delta^n:H\longrightarrow H^{\otimes(n+1)}
\] 为迭代余乘法，设 \(\overline{\Delta}^{\,n}\) 为 \(\Delta^n\) 与投影 \[
  \pi^{\otimes(n+1)}:
  H^{\otimes(n+1)}
  \longrightarrow
  (H^+)^{\otimes(n+1)},
\] 的复合，并置 \[
  H_n'
  =
  \ker(\overline{\Delta}^{\,n})
  =
  \left\{
    x\in H\ \middle|\
    \Delta^n(x)\in
    \sum_{i=0}^n
      H^{\otimes(n-i)}\otimes H_0\otimes H^{\otimes i}
  \right\}.
\]。令 \(\overline{\Delta}^{\,0}=\operatorname{id}_H\)，于是 \(H_0'=H_0\)。直接对 \(n\) 归纳可得 \begin{equation}
  H_n
  \subset
  H_n'
  \subset
  (\overline{I^{n+1}})^\perp.
  \tag{*}
  \label{sga3:VIIB:1.3.6.A:star-prime}
\end{equation}。

至此尚未使用 \(H\) 在 \(k\) 上平坦的假设。现在假设 \(H\) 在 \(k\) 上平坦，因而是投射模。由 \SGARefLink{sga3:VIIB:0.2.2}{0.2.2} 得 \(A^\dagger=H\)。我们将证明 \[
  H_n=(\overline{I^{n+1}})^\perp.
\]。对 \(n=0\) 这显然成立。设已对 \(n<r\) 证明。则 \(H_{r-1}\) 是同态 \[
  H\longrightarrow(\overline{I^r})^\dagger.
\] 的核。由于 \(H^+\) 平坦，同态 \[
  (H/H_{r-1})\otimes H^+
  \longrightarrow
  (\overline{I^r})^\dagger\otimes H^+
\] 为单射。另一方面，假设蕴含 \(I\) 是投射伪紧 \(k\)-模，因为它是 \(A=H^*\) 的直因子。因此由 \SGARefLink{sga3:VIIB:0.3.6}{0.3.6} 得 \[
  \bigl(\overline{I^r}\widehat{\otimes}I\bigr)^\dagger
  \simeq
  (\overline{I^r})^\dagger\otimes I^\dagger
  =
  (\overline{I^r})^\dagger\otimes H^+.
\]。正合列 \[
  \overline{I^r}\widehat{\otimes}I
  \longrightarrow A
  \longrightarrow A/\overline{I^{r+1}}
  \longrightarrow0
\] 经对偶便给出正合列 \begin{equation}
  (\overline{I^r})^\dagger\otimes H^+
  \xleftarrow{\ \delta\ }
  H
  \longleftarrow
  (A/\overline{I^{r+1}})^\dagger
  \longleftarrow0,
  \tag{1}
  \label{sga3:VIIB:1.3.6.A:eq1}
\end{equation}，其中 \(\delta\) 是 \(\Delta_H\) 与投影 \begin{equation}
  H\otimes H
  \longrightarrow
  H\otimes H^+
  \longrightarrow
  (H/H_{r-1})\otimes H^+
  \lhook\joinrel\longrightarrow
  (\overline{I^r})^\dagger\otimes H^+.
  \tag{2}
  \label{sga3:VIIB:1.3.6.A:eq2}
\end{equation} 的复合。对每个 \(u\in H\)，\(\Delta(u)\) 在 \(H\otimes H^+\) 上的投影为 \(\Delta(u)-u\otimes\phi\)。式 \textup{\SGARefLink{sga3:VIIB:1.3.6.A:eq1}{(1)}} 和 \textup{\SGARefLink{sga3:VIIB:1.3.6.A:eq2}{(2)}} 表明，若 \[
  u\in
  (A/\overline{I^{r+1}})^\dagger
  =
  (\overline{I^{r+1}})^\perp,
\]，则 \(\Delta(u)-u\otimes\phi\) 属于 \[
  H\otimes H^+
  \longrightarrow
  (H/H_{r-1})\otimes H^+,
\] 的核，即 \(H_{r-1}\otimes H^+\)。故 \(u\in H_r\)。这就完成了 \textup{\SGARefLink{sga3:vii:paragraph:5-1:equation:i}{(i)}} 和 \textup{\SGARefLink{sga3:vii:paragraph:5-1:equation:ii}{(ii)}} 的证明，同时证明了 \[
  H_n=\ker(\overline{\Delta}^{\,n}).
\]。

现在证明 \textup{\SGARefLink{sga3:VIIB:1.3.6.A:item:iii}{(iii)}}。对 \(i\geq0\) 置 \[
  H_i^+=H_i\cap H^+.
\]。设 \(n\geq1\)。对每个 \(x\in H_n\)，元素 \(\overline x=x-\varepsilon(x)\phi\) 属于 \(H_n^+\)，且 \[
  \Delta(x)
  =
  \varepsilon(x)\phi\otimes\phi
  +\overline x\otimes\phi
  +\phi\otimes\overline x
  +\overline{\Delta}(\overline x).
\]。因此只需证明 \[
  \overline{\Delta}(H_n^+)
  \subset
  \sum_{i=1}^{n-1}H_i^+\otimes H_{n-i}^+.
\]。对 \(i=0,\ldots,n-1\)，映射 \(\overline{\Delta}^{\,n}\) 分解如下：\begin{center}\phantomsection \label{sga3:VIIB:1.3.6.A:iterated-reduced-comultiplication}
% SGA3-VII-DIAG-089
% SGA3-VII-DIAG-089
% Authority: Exp7B-13oct24.pdf, local p. 33.
% Crop: figures/Exp7B_localp033_iterated_reduced_comultiplication_diagram.png
% Locator: tex/components/28_expose_VIIB_section136_filtered_flat_coalgebras_en.tex, Loop 1 line 364.
\begin{tikzcd}[row sep=3.4em,column sep=3.8em]
H^+ \arrow[r,"\overline\Delta"]
  & H^+\otimes H^+
      \arrow[r,"{\overline\Delta^{\,i}\otimes
                    \overline\Delta^{\,n-i-1}}"]
      \arrow[d]
  & (H^+)^{\otimes(i+1)}\otimes(H^+)^{\otimes(n-i)} \\
{}
  & \displaystyle
    \frac{H^+}{H_i^+}\otimes\frac{H^+}{H_{n-i-1}^+}
      \arrow[r,"f"]
  & \displaystyle
    \frac{H^+}{H_i^+}\otimes(H^+)^{\otimes(n-i)}
      \arrow[u,"g"]
      \, .
\end{tikzcd}
\end{center}此外，由于 \(H^+/H_i^+\) 和 \((H^+)^{\otimes(n-i)}\) 平坦，此图中的映射 \(f\) 和 \(g\) 均为单射。因此，对每个 \(i=0,\ldots,n-1\) 都有 \[
  \overline{\Delta}(H_n^+)
  \subset
  H_i^+\otimes H^+
  +
  H^+\otimes H_{n-i-1}^+
\]。在下一个引理中取 \(M=H^+\)、\(E_0=0\)，并对 \(i\geq1\) 取 \(E_i=H_{i-1}^+\)，便得 \textup{\SGARefLink{sga3:VIIB:1.3.6.A:item:iii}{(iii)}}。

\medskip\noindent\textbf{引理 1.3.6.B。}\phantomsection \label{sga3:VIIB:1.3.6.B} {\itshape 设 \(k\) 为环，\[
  0=E_0\subset E_1\subset\cdots\subset E_n\subset M
\] 为 \(k\)-模。假设对每个 \(i\)，\(M/E_i\) 都平坦。则\footnote{法文原文左端第二个加项印作 \(M\otimes E_i\)，而非 \(M\otimes E_{n-i}\)。下面的证明，特别是其中对 \(M/E_{n-i+1}\) 的使用，要求采用后一种下标，此处据此处理。}\[
  \bigcap_{i=0}^n
    \bigl(E_i\otimes M+M\otimes E_{n-i}\bigr)
  =
  \sum_{i=1}^n E_i\otimes E_{n-i+1}.
\]}

\textit{证明。}将左右两端分别记为 \(K\) 和 \(S\)。显然 \(S\subset K\)；我们证明反向包含。对 \(i=0,\ldots,n\) 置 \[
  K_i=K\cap(E_i\otimes M).
\]。对 \(i=1,\ldots,n\)，\(M/E_{n-i+1}\) 和 \(E_i/E_{i-1}\) 均平坦。因此下面的同态 \(\tau_i\) 为单射，而复合映射 \[
  (E_i/E_{i-1})\otimes M
  \longrightarrow
  (E_i/E_{i-1})\otimes(M/E_{n-i+1})
  \xhookrightarrow{\ \tau_i\ }
  (M/E_{i-1})\otimes(M/E_{n-i+1})
\] 的核为 \[
  (E_i/E_{i-1})\otimes E_{n-i+1}.
\]。\(K_i\) 在 \((E_i\otimes M)/(E_{i-1}\otimes M)\) 中的像既包含于此核，又包含此核。故 \[
  K_i=K_{i-1}+E_i\otimes E_{n-i+1},
\]，引理得证。

本段最后回到任意伪紧环 \(k\) 的情形。\footnote{\label{sga3:VIIB:nde:66}下文为编者所增；原文仅处理 \(k\) 为 Artin 环的情形。}设 \((\mathbf H,\Delta,\varepsilon)\) 为平坦 \(\mathbf O_k\)-余代数，置 \(\mathbf H^+=\ker(\varepsilon)\)，并设 \[
  A=\Gamma^*(\mathbf H)
\] 为对偶预有限 \(k\)-代数。置 \(X=\operatorname{Spf}(A)\)，于是 \(\mathbf H=\mathbf H(X)\) \textup{（参见 \SGARefLink{sga3:VIIB:1.3.5}{1.3.5}）}。设给定 \(\mathbf O_k\)-余代数态射 \[
  \phi:\mathbf O_k\longrightarrow\mathbf H
\]。它定义连续的 \(k\)-代数同态 \(A\to k\)，从而定义结构态射 \(X\to\operatorname{Spf}(k)\) 的截面 \(\sigma\)。

对 \(\operatorname{Alf}/k\) 的每个对象 \(B\)，置 \[
  \mathbf H_0(B)=\phi(B)=B\phi_B,
  \qquad
  \phi_B=\phi(1_B)\in\mathbf H(B).
\]。对 \(n\geq1\)，用 \[
  \mathbf H_n(B)
  =
  \left\{
    u\in\mathbf H(B)\ \middle|\
    \Delta(u)-u\otimes\phi_B
    \in
    \mathbf H_{n-1}(B)\otimes\mathbf H^+(B)
  \right\}.
\] 定义 \(\mathbf H\) 的 \(\mathbf O_k\)-子模 \(\mathbf H_n\)。这样得到 \(\mathbf H(X)\) 的滤过 \[
  \mathbf H_0\subset\mathbf H_1\subset\cdots
\]。上述结果给出：

\medskip\noindent\textbf{命题。}\phantomsection \label{sga3:VIIB:1.3.6.proposition} {\itshape 截面 \(\sigma\) 在底拓扑空间之间诱导同构，当且仅当 \[
  \mathbf H(X)=\bigcup_n\mathbf H_n.
\]。}


% 下一游标： Exp7B 分篇页 34, 合并阅读本页 547, 再版印刷页 536, 紧接第 1.4 节之前。
